\sectionheading{5.26 Python Programming Language-এর বৈশিষ্ট্য}

Python জনপ্রিয় হওয়ার অন্যতম কারণ হলো এর
\textbf{সহজ Syntax, পাঠযোগ্য Code, কম Code-এ কাজ করার সুবিধা}
এবং বিভিন্ন ধরনের Application তৈরির উপযোগিতা।

%================================================
% Main Features
%================================================

\booksubheading{Python-এর প্রধান বৈশিষ্ট্য}

%================================================
% Feature 1
%================================================

\textbf{\color{darkTeal}১. সহজ ও বোধগম্য (Simple and Easy)}

Python-এর Syntax তুলনামূলক সহজ এবং সাধারণ English-এর কাছাকাছি।
ফলে নতুনদের জন্য Code পড়া, লেখা ও শেখা সহজ।

\begin{examplebox}{ছোট উদাহরণ}

\begin{codebox}

\begin{english}
\ttfamily
print("Hello")
\end{english}

\end{codebox}

একটি Output প্রদর্শনের জন্য খুব অল্প Code প্রয়োজন।

\end{examplebox}

%================================================
% Feature 2
%================================================

\textbf{\color{darkTeal}২. উচ্চস্তরের ভাষা (High Level Language)}

Python একটি \textbf{High Level Programming Language}।

তাই Programmer-কে সাধারণত Memory Address বা Hardware-এর
নিম্নস্তরের বিষয় সরাসরি নিয়ন্ত্রণ করতে হয় না।

\begin{rememberbox}

\textbf{\color{deepOrange}মনে রাখবে}

High Level
\(\longrightarrow\)
মানুষের জন্য তুলনামূলক সহজে বোঝা ও লেখা যায়।

\end{rememberbox}

%================================================
% Feature 3
%================================================

\textbf{\color{darkTeal}৩. Interpreter-ভিত্তিক ভাষা
(Interpreted Language)}

Python Program সাধারণত
\textbf{Python Interpreter}-এর মাধ্যমে Execute করা হয়।

Program চালানোর সময় Interpreter Code-এর Instructionগুলো
বিশ্লেষণ ও Execute করে।

\begin{rememberbox}

\textbf{\color{deepOrange}MCQ তথ্য}

Python
\(\longrightarrow\)
\textbf{Interpreted Language}

\end{rememberbox}

%================================================
% Feature 4
%================================================

\textbf{\color{darkTeal}৪. General Purpose Language}

Python কোনো একটি নির্দিষ্ট কাজের জন্য সীমাবদ্ধ নয়।

এটি বিভিন্ন ধরনের Program ও Software তৈরিতে ব্যবহার করা যায়।

যেমন—

\begin{itemize}

\item Web Application

\item Automation

\item Data Analysis

\item Artificial Intelligence

\item Machine Learning

\item Scientific Computing

\end{itemize}

\begin{center}

\textbf{একটি ভাষা}
\(\longrightarrow\)
\textbf{বিভিন্ন ধরনের কাজ}

\end{center}

%================================================
% Feature 5
%================================================

\textbf{\color{darkTeal}৫. পাঠযোগ্য Code (Readable)}

Python-এর Code তুলনামূলক পরিষ্কার ও সহজে পড়া যায়।

অন্যান্য অনেক Programming Language-এর তুলনায় একই কাজ
কম ও সরল Code-এর মাধ্যমে প্রকাশ করা সম্ভব।

\begin{center}

\textbf{সহজ Syntax}
\(\longrightarrow\)
\textbf{পাঠযোগ্য Code}

\end{center}

%================================================
% Feature 6
%================================================

\textbf{\color{darkTeal}৬. কম Code-এ Program লেখা যায়}

Python-এ অনেক কাজ তুলনামূলকভাবে অল্প সংখ্যক Statement
ব্যবহার করে করা যায়।

ফলে—

\begin{itemize}

\item Code ছোট হয়।

\item লেখা সহজ হয়।

\item পড়া সহজ হয়।

\item পরিবর্তন ও রক্ষণাবেক্ষণ সহজ হয়।

\end{itemize}

%================================================
% Feature 7
%================================================

\textbf{\color{darkTeal}৭. Case Sensitive}

Python একটি
\textbf{Case Sensitive Programming Language}।

অর্থাৎ বড় হাতের ও ছোট হাতের অক্ষরকে Python ভিন্নভাবে বিবেচনা করে।

\begin{examplebox}{উদাহরণ}

\begin{codebox}

\begin{english}
\ttfamily
name\\
Name\\
NAME
\end{english}

\end{codebox}

উপরের তিনটি নাম Python-এর কাছে এক নয়।

\end{examplebox}

\begin{rememberbox}

\textbf{\color{deepOrange}সতর্কতা}

\textenglish{\ttfamily age}
এবং
\textenglish{\ttfamily Age}

একই Variable নয়।

\medskip

\textbf{Case Sensitive}
\(\longrightarrow\)
Capital ও Small Letter গুরুত্বপূর্ণ।

\end{rememberbox}

%================================================
% Feature 8
%================================================

\textbf{\color{darkTeal}৮. Dynamically Typed}

Python-এ Variable তৈরি করার সময় সাধারণভাবে
আলাদাভাবে Data Type ঘোষণা করতে হয় না।

Variable-এ যে ধরনের Value রাখা হয়, Python সেই অনুযায়ী
তার Type নির্ধারণ করে।

\begin{examplebox}{উদাহরণ}

\begin{codebox}

\begin{english}
\ttfamily
x = 10\\
name = "Rahim"
\end{english}

\end{codebox}

এখানে Variable তৈরির আগে আলাদাভাবে
\textenglish{\ttfamily int}
বা
\textenglish{\ttfamily string}
লিখে Type ঘোষণা করা হয়নি।

\end{examplebox}

\begin{rememberbox}

এই বিষয়টি
\textbf{Variable ও Data Type}
অংশে বিস্তারিতভাবে আলোচনা করা হবে।

\end{rememberbox}

%================================================
% Feature 9
%================================================

\textbf{\color{darkTeal}৯. Object-Oriented Programming সমর্থন করে}

Python
\textbf{Object-Oriented Programming (OOP)}
সমর্থন করে।

অর্থাৎ Class ও Object ব্যবহার করে Program তৈরি করা সম্ভব।

HSC পর্যায়ে এখানে শুধু এই মূল ধারণাটি মনে রাখলেই যথেষ্ট।

%================================================
% Feature 10
%================================================

\textbf{\color{darkTeal}১০. Portable / Cross-platform}

এক Operating System-এ লেখা Python Program প্রয়োজনীয়
Python Environment থাকলে অন্য Operating System-এও
ব্যবহার করা যায়।

\begin{center}

\textbf{Windows}
\(\longleftrightarrow\)
\textbf{Linux}
\(\longleftrightarrow\)
\textbf{macOS}

\end{center}

এ কারণে Python-কে
\textbf{Portable} বা
\textbf{Cross-platform}
Language বলা হয়।

%================================================
% Feature 11
%================================================

\textbf{\color{darkTeal}১১. Free ও Open Source}

Python বিনামূল্যে ব্যবহার করা যায় এবং এর Source Code উন্মুক্ত।

অর্থাৎ—

\begin{itemize}

\item Python ব্যবহার করতে Licence Fee লাগে না।

\item শেখা ও Software Development-এর জন্য সহজে ব্যবহার করা যায়।

\end{itemize}

%================================================
% Feature 12
%================================================

\textbf{\color{darkTeal}১২. বিশাল Library Support}

Python-এর Standard Library এবং বিভিন্ন অতিরিক্ত
Package ও Library রয়েছে।

এসব ব্যবহার করে অনেক সাধারণ ও জটিল কাজ সহজে করা যায়।

যেমন—

\begin{itemize}

\item Mathematical কাজ

\item File Processing

\item Web Development

\item Data Analysis

\item Artificial Intelligence ও Machine Learning

\end{itemize}

\begin{center}

\textbf{Library ব্যবহার}
\(\longrightarrow\)
\textbf{কম Code}
\(\longrightarrow\)
\textbf{দ্রুত Development}

\end{center}

%================================================
% Feature 13
%================================================

\textbf{\color{darkTeal}১৩. Automatic Memory Management}

Python Memory ব্যবস্থাপনার অনেক কাজ স্বয়ংক্রিয়ভাবে করে।

ফলে সাধারণ Programming-এর ক্ষেত্রে Programmer-কে
Memory Allocation ও Deallocation নিয়ে সরাসরি বেশি কাজ
করতে হয় না।

%================================================
% Summary Table
%================================================

\booksubheading{Python-এর বৈশিষ্ট্য এক নজরে}

\begin{tabularx}
{\linewidth}
{
|>{\bfseries\RaggedRight\arraybackslash}p{0.35\linewidth}
|Y|
}

\hline

\rowcolor{softBlue}

\multicolumn{1}{|c|}{\textbf{বৈশিষ্ট্য}}
&
\multicolumn{1}{c|}{\textbf{সহজ অর্থ}}

\\

\hline

Simple and Easy
&
শেখা, পড়া ও লেখা সহজ
\\

\hline

High Level
&
মানুষের জন্য তুলনামূলক সহজে বোঝা যায়
\\

\hline

Interpreted
&
Interpreter-এর মাধ্যমে Program Execute হয়
\\

\hline

General Purpose
&
বিভিন্ন ধরনের কাজে ব্যবহার করা যায়
\\

\hline

Readable
&
Code সহজে পড়া ও বোঝা যায়
\\

\hline

Less Code
&
তুলনামূলক কম Code-এ কাজ করা যায়
\\

\hline

Case Sensitive
&
Capital ও Small Letter-কে ভিন্ন হিসেবে বিবেচনা করে
\\

\hline

Dynamically Typed
&
Variable-এর Type আগে আলাদাভাবে ঘোষণা করা বাধ্যতামূলক নয়
\\

\hline

Object-Oriented
&
Class ও Object ব্যবহার সমর্থন করে
\\

\hline

Portable
&
বিভিন্ন Operating System-এ ব্যবহারযোগ্য
\\

\hline

Free and Open Source
&
বিনামূল্যে ব্যবহারযোগ্য এবং Source Code উন্মুক্ত
\\

\hline

Library Support
&
বিভিন্ন কাজের জন্য প্রচুর Library ও Package রয়েছে
\\

\hline

Automatic Memory Management
&
Memory ব্যবস্থাপনার অনেক কাজ স্বয়ংক্রিয়ভাবে সম্পন্ন হয়
\\

\hline

\end{tabularx}

%================================================
% Important for Exam
%================================================

\begin{rememberbox}

\textbf{\color{deepOrange}পরীক্ষায় বেশি গুরুত্বপূর্ণ}

Python হলো—

\begin{center}

\textbf{High Level}
\(\longrightarrow\)
\textbf{Interpreted}
\(\longrightarrow\)
\textbf{General Purpose}

\medskip

\textbf{Case Sensitive}
\(\longrightarrow\)
\textbf{Dynamically Typed}
\(\longrightarrow\)
\textbf{Object-Oriented}

\medskip

\textbf{Portable}
\(\longrightarrow\)
\textbf{Free and Open Source}

\end{center}

\end{rememberbox}

%================================================
% Common Mistakes
%================================================

\booksubheading{সাধারণ ভুল}

\begin{compacttwocolpoints}

\item Python-কে Low Level Language মনে করা।

\item Python-কে শুধু একটি নির্দিষ্ট কাজের Language মনে করা।

\item Python-কে Case Insensitive মনে করা।

\item
\textenglish{\ttfamily name}
ও
\textenglish{\ttfamily Name}
-কে একই Variable মনে করা।

\item Python-এ Variable-এর Data Type আগে ঘোষণা করতেই হবে
এমন ধারণা করা।

\item Python শুধু Artificial Intelligence-এর জন্য ব্যবহৃত হয়
এমন মনে করা।

\item Portable বলতে শুধু একটি Operating System-এ কাজ করে
এমন মনে করা।

\end{compacttwocolpoints}

%================================================
% Final Reminder
%================================================

\begin{rememberbox}

\textbf{\color{deepOrange}মনে রাখবে}

\begin{itemize}

\item
Python-এর Syntax সহজ ও Code পাঠযোগ্য।

\item
Python একটি
\textbf{High Level Programming Language}।

\item
Python
\textbf{Interpreter}-এর মাধ্যমে Program Execute করে।

\item
Python একটি
\textbf{General Purpose Programming Language}।

\item
Python
\textbf{Case Sensitive}।

\item
Python
\textbf{Dynamically Typed}।

\item
Python
\textbf{Object-Oriented Programming}
সমর্থন করে।

\item
Python
\textbf{Portable / Cross-platform}।

\item
Python
\textbf{Free ও Open Source}।

\item
Python-এর প্রচুর
\textbf{Library ও Package}
রয়েছে।

\item
Python Memory Management-এর অনেক কাজ
স্বয়ংক্রিয়ভাবে করে।

\item
দ্রুত মনে রাখবে—

\begin{center}

\textbf{Simple}
\(\longrightarrow\)
\textbf{High Level}
\(\longrightarrow\)
\textbf{Interpreted}
\(\longrightarrow\)
\textbf{General Purpose}

\medskip

\textbf{Case Sensitive}
\(\longrightarrow\)
\textbf{Dynamic}
\(\longrightarrow\)
\textbf{OOP}
\(\longrightarrow\)
\textbf{Portable}

\end{center}

\end{itemize}

\end{rememberbox}

\small

%================================================
% 5.27 Programming Environment
%================================================

